Reinforcement learning algorithms are asymptotic approximations to classical dynamic programming operators. Value iteration becomes Q-learning when expectations are replaced by single samples; policy iteration becomes the natural policy gradient when the greedy improvement step is approximated by gradient ascent. The stochastic approximation framework guarantees that under appropriate step-size conditions, noisy iterates converge to the same fixed points as their deterministic counterparts. This survey has traced these connections through the theory of contractions and gradient domination, through structural estimation where RL is a computational tool for solving high-dimensional Bellman equations, through strategic games where independent learners converge to equilibrium, and through preference learning where the Bradley-Terry model bridges discrete choice econometrics and RLHF. The practical vulnerabilities of these methods are real. Deep RL lacks convergence guarantees outside tabular settings, results are sensitive to hyperparameters and random seeds, and the absence of standardized simulators for economic environments remains a binding constraint. When guided by economic structure, reinforcement learning extends the reach of dynamic programming to problems that were previously intractable.
